
\section{Refresh Token}

\textit{Normative language per \href{https://tools.ietf.org/html/rfc2119}{RFC~2119}.}

\subsection{Purpose}

Refresh tokens allow the Canvasflow SPA to obtain new access tokens silently --- without
requiring the user to re-authenticate --- after the current access token expires. This
section defines what the tenant's IdP MUST support for refresh token operation, how the
entitlement claim is re-evaluated at refresh time, and where each token type travels in
the system.

\subsection{Refresh Grant Requirements}

The tenant's IdP token endpoint MUST support \texttt{grant\_type=refresh\_token} for the
Canvasflow SPA's public OAuth client registration 
(\href{https://tools.ietf.org/html/rfc6749\#section-6}{RFC~6749~\S6}).

\subsubsection{Refresh Request}

The SPA sends the following request to the tenant's token endpoint when refreshing:

\begin{lstlisting}
POST <tenant token endpoint>
Content-Type: application/x-www-form-urlencoded

grant_type=refresh_token
&refresh_token=<stored_refresh_token>
&client_id=<canvasflow_client_id>
&scope=openid
\end{lstlisting}

No \texttt{client\_secret} is included --- the client is a public client (\hyperref[sec:token-endpoint-authentication]{\S7.6}).

\subsubsection{Refresh Response}

The token endpoint MUST respond with a JSON body containing at minimum:

\begin{lstlisting}[language=json]
{
  "access_token": "<new_jwt_access_token>",
  "token_type": "Bearer",
  "expires_in": 3600,
  "refresh_token": "<new_or_same_refresh_token>"
}
\end{lstlisting}

The \texttt{access\_token} in the response MUST be a conforming JWT access token
satisfying all requirements in \hyperref[sec:token-requirements]{\S5 Token Requirements} --- including the
\texttt{cf:entitlements} (or equivalent) claim with freshly evaluated entitlement values.

\subsection{Freshly Evaluated Entitlements on Every Refresh}\label{sec:freshly-evaluated-entitlements}

The access token issued in response to a refresh grant MUST carry an entitlement claim
that was freshly evaluated at refresh time. The pre-token hook MUST re-run on every
refresh and MUST read from the tenant's live source of truth. The hook MUST NOT copy the
entitlement claim from the previous access token.

\textbf{Why this is the revocation mechanism.} The Canvasflow Resource Server has no
revocation endpoint and no token blocklist (\hyperref[sec:out-of-scope]{\S1.4}, \hyperref[sec:canvasflow-responsibilities]{\S3.3}). When a tenant removes a
user's entitlement in their system, the change takes effect at the user's next refresh
--- when the hook re-evaluates and issues a new access token without the removed
entitlement. If the tenant also revokes the user's refresh token in the IdP, the user
cannot obtain any further access tokens, and the residual exposure window is bounded by
the remaining lifetime of the current access token (\texttt{exp}, $\leq$ 3600~seconds).
This is the complete revocation workflow available in this version.

Failure to re-evaluate on refresh means a user whose access has been revoked in the
tenant's system will continue to receive premium content for as long as their refresh
token remains valid --- potentially indefinitely. The hook re-evaluation requirement is
therefore a security requirement, not only an operational one.

\subsection{Refresh Token Rotation}

Refresh token rotation is RECOMMENDED but not required.

When the IdP rotates refresh tokens (issues a new refresh token on each refresh grant
and invalidates the old one), the SPA MUST store the new refresh token received in the
refresh response, replacing the previous value in its storage.

If the IdP implements refresh token rotation with an absolute maximum lifetime or a
sliding-window reuse detection policy, the tenant MUST document this behavior for
Canvasflow support, as it affects how silent session renewal is planned.

\subsection{Token Routing Table}

Each of the three tokens issued by the tenant's IdP has a defined routing scope. Tokens
MUST NOT be sent to endpoints outside their defined routing scope.

\begin{table}[htbp]
\centering
\begin{tabularx}{\textwidth}{@{}lXXX@{}}
\toprule
\textbf{Token} & \textbf{Purpose} & \textbf{MUST be sent to} & \textbf{MUST NOT be sent to} \\
\midrule
\textbf{Access token} &
  Authorize API requests &
  Canvasflow GraphQL API (as \texttt{Authorization: Bearer}) &
  Tenant IdP; any other endpoint \\[6pt]
\textbf{ID token} &
  Prove identity to the SPA for display purposes &
  SPA only (for reading user display name, picture, etc.) &
  Canvasflow GraphQL API; tenant IdP token endpoint \\[6pt]
\textbf{Refresh token} &
  Obtain new access tokens &
  Tenant IdP token endpoint (\scriptsize \texttt{grant\_type=refresh\_token}) &
  Canvasflow GraphQL API; any other endpoint \\
\bottomrule
\end{tabularx}
\caption{Token routing: where each token travels}
\end{table}

The Canvasflow GraphQL API MUST NOT receive the ID token or the refresh token under any
circumstances. An ID token presented to the API would fail the \texttt{typ: "at+jwt"}
check and be rejected with 401. A refresh token presented to the API would similarly be rejected.

\subsection{Client Storage (Informative)}

The following describes the Canvasflow PWA's token storage behavior. This section is
informative --- it describes Canvasflow's client-side implementation for tenant
understanding. It does not impose requirements on the tenant.

\begin{itemize}
  \item The \textbf{access token} is held in memory only (a JavaScript variable). It is
    not written to \texttt{localStorage}, \texttt{sessionStorage}, cookie storage, or
    any other durable browser storage. If the page is closed or refreshed, the access
    token is lost and the SPA uses the refresh token to obtain a new one.

  \item The \textbf{refresh token} is persisted to secure, same-origin storage (the
    specific mechanism is an implementation detail of the Canvasflow PWA). It survives
    page navigation and browser restarts.

  \item The \textbf{ID token} is used briefly after the token exchange to populate user
    display information and then discarded. It is not persisted.
\end{itemize}

This pattern mirrors the storage recommendations of OAuth~2.0 for Browser-Based Apps
(\href{https://tools.ietf.org/html/rfc9449}{RFC~9449} and the draft OAuth~2.0 for Native Apps guidelines), where long-lived
credentials are handled in memory and short-lived credentials are persisted only as
needed for session continuity.